\chapter{Complex numbers: multiplication, length, and conjugation}
\label{learn:complex_numbers}
\begin{learningcard}{Bring forward and choose your next operation}
\textsf{\textbf{Question.}} How does multiplication change a complex number?\par
\textbf{Required lessons.} \hyperref[learn:proof]{From convincing examples to proof}\par
\textbf{Helpful background.} A coordinate-plane sketch helps interpret multiplication.\par
\textbf{Learn to.} Compute conjugates, lengths, division, and normalized complex columns.\par
\textbf{Entry check.} Expand (a+b)(c+d).\par
\hyperref[solution:entry:complex_numbers]{Check your answer}; repair the operation in \href{https://github.com/Oichkatzelesfrettschen/precalc_paper}{Companion: Precalculus}.
\end{learningcard}


\section{How can arithmetic beyond real roots retain a positive length?}
The companion's completing-the-square lesson introduces $i^2=-1$ to solve
an equation such as $x^2+4=0$. Carry that arithmetic forward: a complex
number $z=a+bi$ has real coordinates $a,b$. Equality means equality of
both coordinates. A real number embeds as $a+0i$.
Our question now concerns measurement. Squaring $i$ gives $-1$, yet a
squared length should be nonnegative. Which product measures length?

\section{Distribute, then replace the square of \texorpdfstring{$i$}{i}}
Add and subtract corresponding real and imaginary parts. For multiplication,
use distribution before collecting terms:
\[
 (a+bi)(c+di)=ac+adi+bci+bdi^2=(ac-bd)+(ad+bc)i.
\]
\lessoncue{Worked example}{Separate the arithmetic operations.}
For $z=2+i$ and $w=1-3i$, addition gives $z+w=3-2i$.
Multiplication gives
\[
 zw=2-6i+i-3i^2=5-5i.
\]
The second equality replaces $i^2$ by $-1$ and combines imaginary terms.
Multiplying a general number by $i$ gives $i(a+bi)=-b+ai$: its plane
coordinates change from $(a,b)$ to $(-b,a)$, a quarter-turn counterclockwise.

\begin{figure}[htbp]
\centering
\begin{tikzpicture}[x=1.25cm,y=1.25cm,>=Stealth,font=\small]
\draw[->] (-1.8,0)--(3,0) node[right] {real};
\draw[->] (0,-1.6)--(0,2.7) node[above] {imaginary};
\draw[->,thick] (0,0)--(2,1) node[right] {$z=2+i$};
\draw[->,thick,dashed] (0,0)--(-1,2) node[left] {$iz=-1+2i$};
\draw[->,thick,dotted] (0,0)--(2,-1) node[right] {$\overline z=2-i$};
\draw[->] (0.55,0.275) arc[start angle=26.565,end angle=116.565,radius=0.615];
\node at (-0.05,0.85) {$\times i$};
\foreach \position in {-1,1,2} {
\draw (\position,0.05)--(\position,-0.05) node[below] {$\position$};
\draw (0.05,\position)--(-0.05,\position);
}
\end{tikzpicture}

\caption{The solid arrow represents $z=2+i$. Multiplication by $i$ gives
the dashed arrow $iz=-1+2i$, a quarter-turn of equal length. Reflection
across the real axis gives the dotted arrow $\overline z=2-i$. Line styles
and endpoint labels distinguish the three operations in grayscale.}
\end{figure}

\section{Conjugate to measure length and divide}
The conjugate of $z=a+bi$ is $\overline z=a-bi$. Multiply the pair:
\[
 \overline z z=(a-bi)(a+bi)=a^2+b^2.
\]
The cross terms cancel. Define the modulus $|z|=\sqrt{a^2+b^2}$,
the distance from the origin in the complex plane. Therefore
$|z|^2=\overline z z$, and $|z|=0$ exactly when $z=0$.
For $z=2+i$, $|z|=\sqrt5$, while $z^2=3+4i$ is a different quantity.

\lessoncue{Worked example}{Make a denominator real before dividing.}
For $w\ne0$, multiply numerator and denominator by $\overline w$.
The denominator becomes the positive real number $|w|^2$:
\[
 \frac{2+i}{1-3i}
 =\frac{(2+i)(1+3i)}{(1-3i)(1+3i)}
 =\frac{-1+7i}{10}=-\frac1{10}+\frac7{10}i.
\]
Check by multiplying the answer by $1-3i$:
$(-1+7i)(1-3i)/10=(20+10i)/10=2+i$.
Distribution also verifies $\overline{zw}=\overline z\,\overline w$.
Consequently $|zw|^2=|z|^2|w|^2$, so $|zw|=|z||w|$ because
moduli are nonnegative. Multiplication by $i$ preserves length since $|i|=1$.

\section{Extend the dot product with conjugation}
A complex coordinate vector is a column whose entries are complex numbers.
Add columns component by component and scale each entry as before. For
equal-length columns define
\[
 \langle u,v\rangle=\sum_{s=1}^n\overline{u_s}v_s,
 \qquad \|v\|=\sqrt{\langle v,v\rangle}
 =\sqrt{\sum_{s=1}^n |v_s|^2}.
\]
We conjugate the first argument, the convention used in the quantum bridge.
Thus $\langle u,cv\rangle=c\langle u,v\rangle$, while
$\langle cu,v\rangle=\overline c\langle u,v\rangle$.
Exchanging the arguments conjugates the result:
$\langle v,u\rangle=\overline{\langle u,v\rangle}$.
The conjugate transpose $u^*$ means transpose the column and conjugate
its entries, so $\langle u,v\rangle=u^*v$; $u^\dagger$ is another common
notation. For real columns the formula reduces to the ordinary dot product.

\lessoncue{Worked example}{Test two complex directions.}
Take $u=(1,i)^T$ and $v=(i,1)^T$. Then
\[
 \langle u,v\rangle=1\cdot i+(-i)\cdot1=0,\qquad
 \|u\|^2=1+(-i)i=2.
\]
The columns are orthogonal, meaning their inner product is zero.
Dividing either column by $\sqrt2$ gives a unit vector. Omitting conjugation
would give $u^Tu=1+i^2=0$ for a nonzero column, so that expression fails
as a squared norm. In quantum mechanics the normalized column
$u/\sqrt2$ has squared component moduli $1/2,1/2$. Their interpretation as
measurement probabilities requires the state and measurement assumptions
introduced in Chapter~\ref{learn:quantum}.

\section{Try and check}
\paragraph{1. Multiply and check the length rule.}
For $z=3-2i$ and $w=-1+i$, find $zw$ and verify $|zw|^2=|z|^2|w|^2$.
\paragraph{Solution.}
Distribution gives $zw=-3+3i+2i-2i^2=-1+5i$.
Its squared modulus is $1+25=26$. The separate squared moduli are
$9+4=13$ and $1+1=2$, whose product is twenty-six.

\paragraph{2. Divide and reconstruct the numerator.}
Compute $(3+i)/(1+i)$.
\paragraph{Solution.}
Multiply both numerator and denominator by $1-i$:
$(3+i)(1-i)/2=(4-2i)/2=2-i$.
Multiplication checks $(2-i)(1+i)=3+i$.

\paragraph{3. Normalize a complex column.}
For $v=(1+i,2)^T$, find a unit vector and verify its squared norm.
\paragraph{Solution.}
The squared norm is $|1+i|^2+|2|^2=2+4=6$.
Set $q=v/\sqrt6$. Its squared norm is $2/6+4/6=1$.
Each component uses its modulus squared; squaring the complex component
itself would give $(1+i)^2=2i$ and lose the length interpretation.

\section*{Ready for the next question?}
\label{readiness:complex_numbers}
Find the conjugate and squared modulus of 3+4i, then its reciprocal.
\par\hyperref[solution:ready:complex_numbers]{Read the worked readiness solution.} Continue when you can explain the operations and assumptions. Otherwise return to the method and its solved exercises.

\lessonreference
{Distribute products using $i^2=-1$; use $\overline z z=|z|^2$ for
length. Divide by $w\ne0$ by multiplying by $\overline w/\overline w$.}
{The complex inner product conjugates its first argument here.
Normalization divides a nonzero vector by its positive real norm.}
{Chapter~\ref{learn:quantum} uses complex amplitudes and inner products;
Chapter~\ref{learn:geometry} uses lengths and products as geometric data.}

